\begin{table}[t]
\centering
\caption{SCL Feature Extraction Front-End: 10-D Time-Domain Features.
All features computed in Python (float64) and SCL (REAL/float32) produce
identical results within IEEE~754 roundtrip ($<10^{-4}$).
End-to-end chain using only 10-D SCL-computed features yields 45.20% accuracy vs.\ 100.00% for
full 28-D Python features. The 54.8%
drop quantifies the cost of omitting frequency-domain and dispersion-entropy
features on the PLC.}
\label{tab:scl_feature_frontend}
\small
\begin{tabular}{@{}lrr@{}}
\toprule
\textbf{Metric} & \textbf{Value} \\
\midrule
Feature extraction (Python f64 vs SCL REAL f32) & \\
~~Max roundtrip error & 0.0e+00 \\
~~Equivalence & IDENTICAL \\
\midrule
End-to-end (10-D SCL features $+$ KAN inference) & \\
~~FP32 accuracy & 0.4440 \\
~~LUT accuracy & 0.4520 \\
~~FP32-LUT agreement & 0.9880 \\
\midrule
Ablation comparison & \\
~~Full 28-D (Python) & 1.0000 \\
~~10-D time-only (PLC) & 0.4520 \\
~~Accuracy drop & 0.5480 \\
\bottomrule
\end{tabular}
\vspace{2pt}
{\scriptsize SCL block: \texttt{FeatureExtraction\_TD} (FUNCTION\_BLOCK, 10-D
time-domain), with \texttt{NeuroPLC\_Inference} (FB2, KAN $[28,16,4]$).
The 10-D features are z-score normalized and zero-padded to 28-D.
Frequency-domain and dispersion-entropy front-ends require FFT (future work).}
\end{table}